\section{Related Work and Novelty Boundaries}
\label{sec:related}
\subsection{Resilience and Application Semantics}
Sterbenz et al.~\cite{sterbenz2010} define network resilience in terms of maintaining acceptable service despite faults and challenges, and distinguish operational network state from service state. Their multilevel framework and service-state analysis predate this benchmark. The minimum-useful vocabulary here is a small executable instance of an established service-sensitive idea, rather than a claim to have originated that idea. The contribution lies in exposing what this particular corpus exercises and in reconciling its numerical reporting with the code.

Edge computing can provide nearby processing and a fallback when remote cloud service is unavailable~\cite{satyanarayanan2017}. This supports the architectural motivation for local continuity, but local execution does not automatically preserve every workload. A local copy must contain the needed content, application, credentials, and adequate resources. The stored profiles encode only a declared coverage flag; they do not establish those conditions through physical execution.

\subsection{Benchmarking and Physical Anchoring}
EdgeDroid~\cite{olguin2019} records sensory traces, replays them through application task models, and collects pipeline timing from real client and cloudlet systems. Its workload-conditioned benchmarking is a close methodological comparator. The present corpus uses labeled snapshots and Bernoulli outcomes, without trace playback or execution of application code. It therefore cannot replace that kind of application measurement. A prospective extension could associate each continuity label with a trace-defined task, while preserving unsupported workloads as failures rather than changing the workload after the outage.

Fuller et al.~\cite{fuller2020} discuss digital-twin definitions, data requirements, and common confusion between twins and other digital representations. Contemporary work by Li et al.~\cite{li2025} proposes a native network digital-twin architecture and reports a comparison with real network testing for a signaling-storm use case. The abstract and official publication metadata of that article were inspected; its detailed experimental methodology was not available for full-text audit in this package. Its reported physical comparison nevertheless illustrates the evidence boundary absent here. We do not reproduce or compare numerically against its system.

\subsection{Claimed Increment}
Table~\ref{tab:novelty} positions the evidence relative to the inspected literature. The novelty claim is intentionally modest: a reproducible profile corpus coupled with an explicit audit of synthetic endpoint semantics. Four-state classification, offline fallback, and digital-twin evaluation are established concepts. No exhaustive priority claim is made. The current study may remain too narrow for a research journal unless reviewers judge the semantic failures and reproducibility contribution sufficiently general or additional validation establishes that generality.

\begin{table*}[t]
\centering\caption{Contribution boundaries against closely related primary literature. These are methodological comparisons, not executed performance baselines.}\label{tab:novelty}
\small\begin{tabular}{p{0.19\textwidth}p{0.34\textwidth}p{0.38\textwidth}}\toprule
Source & Established contribution & Increment and boundary of this study \\\midrule
Sterbenz et al.~\cite{sterbenz2010} & Acceptable service under challenges; operational and service states & Executable small corpus and metric audit; no new general theory of resilience \\
Satyanarayanan~\cite{satyanarayanan2017} & Edge proximity and local fallback during cloud disruption & Workload support flags expose failure; no measured edge or cloud comparison \\
EdgeDroid~\cite{olguin2019} & Trace replay, task models, and measured application timing & Snapshot semantics and synthetic occupancy; no claim of equivalent workload realism \\
Fuller et al.~\cite{fuller2020} & Twin definitions, enabling technologies, and data-quality challenges & Clear physical-calibration boundary; no sensor-synchronized twin \\
Li et al.~\cite{li2025} & Native network twin with reported real-network comparison & Reproducibility audit only; no equivalent network-validation evidence \\
\bottomrule\end{tabular}
\end{table*}
